% Source: https://artofproblemsolving.com/community/c3h1987179_geometry_construction_problem
% Stable identity: aops-topic-1987179-q01
\textbf{No.}

Suppose a regular heptagon $V_0V_1\ldots V_6$ could be constructed in a circle; scale so that the circle has center $O$ and radius $1$. Dropping the perpendicular from $V_1$ to $OV_0$ (a ruler-and-compass construction) would produce a segment of length $\cos\dfrac{2\pi}{7}$, so $c=2\cos\dfrac{2\pi}{7}$ would be constructible.

\begin{asy}
usepackage("amsmath");
size(4.5cm);
defaultpen(fontsize(10pt));
draw(unitcircle,gray(0.6));
pair[] v; for(int k=0;k<7;++k) v.push(dir(360*k/7));
draw(v[0]--v[1]--v[2]--v[3]--v[4]--v[5]--v[6]--cycle,linewidth(1));
for(int k=0;k<7;++k) dot(v[k]);
dot((0,0)); label("$O$",(0,0),W);
draw((0,0)--v[0],dashed); draw((0,0)--v[1],dashed);
pair Fv=(v[1].x,0); draw(v[1]--Fv,dotted); dot(Fv);
draw((0,0)--Fv,red+linewidth(1.1));
label("$\cos\dfrac{2\pi}{7}$",Fv/2,S,red+fontsize(8pt));
label("$V_0$",v[0],E); label("$V_1$",v[1],NE);
label("\mbox{}",(min(currentpicture,user=true)+max(currentpicture,user=true))/2,black); // keeps the text after the figure black when LaTeX typesets the labels inline
\end{asy}

\emph{An equation for $c$.} Let $\zeta=\cos\dfrac{2\pi}{7}+i\sin\dfrac{2\pi}{7}$. Then $\zeta^7=1$ and $\zeta\ne1$, so $1+\zeta+\zeta^2+\dots+\zeta^6=0$. Dividing by $\zeta^3$,
\[
\left(\zeta^3+\zeta^{-3}\right)+\left(\zeta^2+\zeta^{-2}\right)+\left(\zeta+\zeta^{-1}\right)+1=0 .
\]
With $c=\zeta+\zeta^{-1}$ we have $\zeta^2+\zeta^{-2}=c^2-2$ and $\zeta^3+\zeta^{-3}=c^3-3c$, so
\[
c^3+c^2-2c-1=0 .
\]
\emph{No rational root.} A rational root of this monic integer cubic would be an integer dividing $1$, but $1+1-2-1=-1\ne0$ and $-1+1+2-1=1\ne0$.

By the quoted theorem (Wantzel), no root of this cubic is constructible. This contradicts the constructibility of $c$, so the regular heptagon cannot be constructed with ruler and compass. $\blacksquare$ (By the Gauss--Wantzel theorem, a regular $n$-gon is constructible exactly when $n$ is a power of $2$ times a product of distinct Fermat primes $3,5,17,257,65537$; $7$ is not a Fermat prime.)
